%==============================================================================
%  sections_v3/proofs_theorem_ge.tex -- proofs of the paper's Theorem (T1) and
%  of the general-equilibrium Proposition (C1).
%
%  SOURCE OF TRUTH: research/model_v4/MODEL_CARD.md, version stamp 2026-08-30
%  (F5 route R40-A, gate PASS, theory-record freeze;
%  commit <pending-orchestrator-hash>), section 6 result ledger.  The two
%  proofs below establish exactly the statements carried in
%  sections_v3/theorem_section.tex; no statement is strengthened, weakened or
%  re-scoped here.  Sources: research/model_v4/proofs/T1_proof.md (fix round
%  closed, items N1-N4 applied) and research/model_v4/proofs/C1_proof.md with
%  research/model_v4/rederive/C1_rederivation.md (additions N1, N2).
%
%  Internal labels are prefixed eq:T1- and eq:C1-.  No environment is defined
%  here; amsthm's proof environment only.
%==============================================================================

\begin{proof}[Proof of Theorem~\ref{thm:T1} (T1)]
Throughout, policies are frozen in the sense of hypothesis~(T-1): the plan menu $\mathcal J$, the
execution policies $B_j(\cdot,\cdot)$, $b_j^*(\cdot)$ and $Q_j^F(\cdot)$, and the cutoff vector $k$
are held fixed in $\kappa$ and held fixed across the two rules compared at each margin. Nothing
below lets $k$ re-solve $k = \Tmap(k;\vartheta)$ at the second rule; that is the whole difference
between this theorem and Proposition~\ref{prop:C1}.

Two pieces of hypothesis traffic are worth naming before the argument starts, because several steps
draw on them. First, hypothesis~(T-5) imports Lemma~\ref{lem:L2} \emph{with its own hypotheses
travelling}, and those hypotheses are Assumption~\ref{asm:A1}, Assumption~\ref{asm:A2p} together
with the primitive table restrictions of Assumption~\ref{asm:TR}, Assumption~\ref{asm:A4},
Assumption~\ref{asm:A5}, Assumption~\ref{asm:A7p} in its on-path injective form,
Lemma~\ref{lem:D1}, the no-feedback timing of Assumption~\ref{asm:nofeedback}, $\Omega > 0$, and an
explicit bidder-entry rule. Where a step below consumes one of them it is cited by name, and it is
consumed as a travelling hypothesis of~(T-5) rather than as a free-standing assumption of this
theorem. The A7 form matters and is therefore named every time: the form
Lemma~\ref{lem:L2} consumes is Assumption~\ref{asm:A7p} (A7$'$), not the weak identification
wording of Assumption~\ref{asm:A7}, which permits two pairs $(j,s)$ with different pooled paths and
is that lemma's first failure case. Second, hypothesis~(T-8) --- $\kappa$-differentiability of
$M_P$ --- is supplied by no standing hypothesis of Section~\ref{sec:hypotheses} and is carried here
as a hypothesis of the theorem; Step~4 is where it is used, and the integrability clause of
Assumption~\ref{asm:A2p} is \emph{not} cited for it, since boundedness is not differentiability.

\medskip
\noindent\emph{Part~(A): the factorisation.}

\emph{Step 1 (the two-cell identity).} By hypothesis~(T-2), Assumption~\ref{asm:A8} holds at the
policy under consideration, so $0 < \Omega < 1$ and Lemma~\ref{lem:L1}, imported by
hypothesis~(T-4), applies in its non-degenerate branch: identity~\eqref{eq:L1} holds at
$(\kappa,\tau,T)$, with $M_F$ and $M_P$ the cell averages of Definition~\ref{def:premium}.

\emph{Step 2 (the flagged cell does not move with $\kappa$).} By Lemma~\ref{lem:L2}, imported with
its hypotheses by~(T-5), and at the fixed cutoff and execution policies of~(T-1), the flagged
posterior, the flagged price, the entry probability and $M_F$ are invariant to $\kappa$. Hence
$\partial_\kappa M_F = 0$. This is the step at which Lemma~\ref{lem:L2}'s stack is consumed, and in
particular Assumption~\ref{asm:A7p}: on a menu that violates it the flagged tuple no longer
identifies the informed component of the selected plan, $M_F$ retains direct $\kappa$-dependence,
and everything after this step is void.

\emph{Step 3 (the flagged weight does not move with $\kappa$ either --- derived, not assumed).}
Hypothesis~(T-7) is PE-$\Omega$, the statement that $\Omega(\cdot,\tau,T)$ does not move with
$\kappa$ at fixed policies. It is derivable at fixed policies, and deriving it is what makes the
partial-equilibrium restriction visible rather than buried. By the clock equivalence of
Lemma~\ref{lem:D1}\,(b), imported by~(T-6), a Voice plan $j$ files by the horizon exactly when
$B_j(s,H-T) \ge \tau$, so with $j(s)$ the plan the frozen cutoff vector assigns to the signal,
\[
D \;=\; \ind\bigl\{a_{j(s)} = 1,\ B_{j(s)}(s,H-T) \ge \tau\bigr\}.
\]
The no-feedback timing of Assumption~\ref{asm:nofeedback}, carried as hypothesis~(T-13), makes
$B_j(s,d)$ a function of $(j,s,d)$ alone: no realised order flow and no realised price enters the
executed path. With the policies frozen by~(T-1), $D$ is therefore a function of $s$ alone. The law
of $s$ carries no $\kappa$: Assumption~\ref{asm:A1} makes $v$ and $\varepsilon$ independent with
strictly positive variances, and the distributional forms of clause~(TR-i) of
Assumption~\ref{asm:TR} --- both travelling with Lemma~\ref{lem:L2} under~(T-5) --- give
$s = v + \varepsilon \sim N(\mu_v,\sigma_v^2+\sigma_\varepsilon^2)$, in which $\kappa$ does not
appear. Indeed $\kappa$ enters the primitives of Definition~\ref{def:primitives} in exactly one
place, the ternary noise-mark law of $z_d$, which reaches observed pooled order flow
$X_d = q_{jd} + z_d$ and reaches nothing that $D$ depends on. Hence $\Omega = \Prb(D=1)$ is
literally the same number at every $\kappa$. The constancy is the form Step~6 needs; the derivative
form $\partial_\kappa\Omega = 0$ is its corollary, and it is the form Step~4 needs.

What this derivation consumes is the freezing of the cutoff vector, not any property of the
disclosure rule. In equilibrium $k$ solves $k = \Tmap(k;\vartheta)$ and moves with $\kappa$,
$\Omega$ moves with it, and the term discarded in Step~4 returns at full size. That term is the
object Proposition~\ref{prop:C1} bounds.

\emph{Step 4 (differentiating in $\kappa$).} Two of the three factors on the right of \eqref{eq:L1}
are constant in $\kappa$: $M_F$ by Step~2 and $\Omega$ by Step~3. Neither step uses the present one,
so nothing here is circular. With those two factors constant, $\kappa \mapsto \Dact$ is an affine
image of $\kappa \mapsto M_P$, which hypothesis~(T-8) makes differentiable; hence
$\partial_\kappa\Dact$ exists and equals $(1-\Omega)\,\partial_\kappa M_P$. Written as one term per
factor that could carry $\kappa$, the same computation reads
\begin{equation}
\label{eq:T1-three}
\partial_\kappa\Dact
\;=\; \Omega\,\partial_\kappa M_F
\;+\; (1-\Omega)\,\partial_\kappa M_P
\;+\; (\partial_\kappa\Omega)\,(M_F - M_P),
\end{equation}
the three terms being the flagged cell's own motion, the pooled cell's own motion, and the
reallocation of mass between the cells. Step~2 kills the first and Step~3 the third. Note what the
third deletion does \emph{not} assume: $M_F - M_P$ is neither small nor signed here, and the term
goes because its coefficient is zero.

\emph{Step 5 (Part~(A)'s factorisation).} By~(T-2), $1-\Omega \in (0,1)$, so
$|1-\Omega| = 1-\Omega$. Taking absolute values in the conclusion of Step~4, and writing
$\mathcal S = |\partial_\kappa\Dact|$ and $\mathcal S_P = |\partial_\kappa M_P|$ as in
Definition~\ref{def:premium},
\begin{equation}
\label{eq:T1-factor}
\mathcal S(\kappa,\tau,T) \;=\; \bigl(1-\Omega(\tau,T)\bigr)\,\mathcal S_P(\kappa,\tau,T),
\end{equation}
exactly, which is the first clause of Part~(A). Step~3 licenses writing the argument of $\Omega$ as
$(\tau,T)$ rather than $(\kappa,\tau,T)$, and that convention is kept from here on. Two consequences
are used later. First, $\mathcal S > 0$ exactly when $\mathcal S_P > 0$, which hypothesis~(T-3)
supplies. Second, when $\mathcal S_P > 0$ the map $x \mapsto |x|$ is differentiable at
$\partial_\kappa M_P \neq 0$, so $\mathcal S_P$ inherits whatever differentiability
$\partial_\kappa M_P$ has in a policy coordinate. That transfer creates no differentiability; it
only moves it, which is why Part~(C)'s local form carries its own smoothness hypothesis.

\emph{Step 6 (the aggregate form, with no differentiability used).} Fix any grid
$\kappa_0 < \kappa_1 < \dots < \kappa_n$ inside the maintained parameter set, with~(T-2) holding at
each node. Applying Step~1 at $\kappa_i$ and at $\kappa_{i+1}$, and using Step~2 and Step~3 in their
\emph{constancy} form --- $M_F$ and $\Omega$ are one and the same number at the two nodes, which a
vanishing derivative at a single $\kappa$ would not deliver --- gives
$\Dact(\kappa_{i+1}) - \Dact(\kappa_i) = (1-\Omega)\bigl[M_P(\kappa_{i+1}) - M_P(\kappa_i)\bigr]$.
Summing absolute values,
\begin{equation}
\label{eq:T1-tv}
\sum_{i=0}^{n-1}\bigl\lvert \Dact(\kappa_{i+1}) - \Dact(\kappa_i)\bigr\rvert
\;=\; (1-\Omega)\sum_{i=0}^{n-1}\bigl\lvert M_P(\kappa_{i+1}) - M_P(\kappa_i)\bigr\rvert,
\end{equation}
which is the total-variation clause of Part~(A). Hypothesis~(T-8) is not used: \eqref{eq:T1-tv} is
finite differences throughout. The same argument runs for any aggregator of the increment vector
that is positively homogeneous of degree one --- total variation, mean absolute slope, the supremum
of the absolute increments --- because $1-\Omega$ is one nonnegative scalar common to every
increment. Every ratio statement in Parts~(B) and~(C) therefore holds verbatim with $\mathcal S$ and
$\mathcal S_P$ replaced by their aggregate counterparts. This matters because a measurement
convention that reports $\kappa$-sensitivity as a total variation over a grid rather than as a
pointwise derivative is then measuring the same functional the theorem is about. Part~(A) is
Steps~5 and~6.

\medskip
\noindent\emph{Part~(B): the threshold margin.}

Throughout Part~(B) the window $T$ is common to the two rules and the comparison is
$b_0 < \tau' < \tau$, so $\tau'$ is the tighter threshold. The maintained configuration
$b_0 < \tau$ is clause~(TR-i) of Assumption~\ref{asm:TR}, and it is imposed at both thresholds.

\emph{Step 7 (the chord form of the pooled sensitivity).} By hypothesis~(T-9),
Assumption~\ref{asm:Atau} holds at each compared policy, with $\bar\pi$ read as the upper support
point of the pooled engagement posterior in \eqref{eq:atau} and not as the pooled engagement share.
Clause~(br-ii) of Assumption~\ref{asm:Abr}, imported by hypothesis~(T-11), places all
$\kappa$-dependence of $M_P$ in the weights of \eqref{eq:atau} and delivers
$\partial_\kappa M_P = \Delta_m\,\Akap\,C_h(\bar\pi)$ exactly, with no
composition-through-$\kappa$ remainder; this is the $M_P$-level form of part~(a) of
Lemma~\ref{lem:L3}, imported by hypothesis~(T-10), which gives
$\partial_\kappa\Eop_\kappa[h] = \Akap\,C_h(\bar\pi)$. Taking absolute values, and using
$\Delta_m > 0$ from clause~(TR-i) of Assumption~\ref{asm:TR} as it travels with~(T-5),
\begin{equation}
\label{eq:T1-chord}
\mathcal S_P \;=\; \Delta_m\,\lvert\Akap\rvert\,\lvert C_h(\bar\pi)\rvert
\;=\; \frac{\Delta_m}{4}\,\lvert\Akap\rvert\,\bar\pi^2\,\lvert h''(\zeta)\rvert
\qquad\text{for some }\zeta\in(0,\bar\pi),
\end{equation}
the second equality by part~(b) of Lemma~\ref{lem:L3}, which is an identity rather than an
approximation. The reading used below is that $\mathcal S_P$ is a \emph{product} of a
weight-derivative magnitude and a chord magnitude. That is why leg~3 of Lemma~\ref{lem:L4} needs
both a coefficient comparison and a chord comparison, and neither alone.

\emph{Step 8 (the ratio identity).} Apply \eqref{eq:T1-factor} at $(\tau',T)$ and at $(\tau,T)$,
which~(T-2) permits at both. By~(T-3), $\mathcal S_P(\kappa,\tau,T) > 0$, so
$\mathcal S(\kappa,\tau,T) > 0$ by \eqref{eq:T1-factor} and the quotient is defined:
\begin{equation}
\label{eq:T1-ratio-tau}
\frac{\mathcal S(\kappa,\tau',T)}{\mathcal S(\kappa,\tau,T)}
\;=\; \frac{\bigl(1-\Omega(\tau',T)\bigr)\mathcal S_P(\kappa,\tau',T)}
            {\bigl(1-\Omega(\tau,T)\bigr)\mathcal S_P(\kappa,\tau,T)}
\;=\; W_\tau\,C_\tau,
\qquad
W_\tau = \frac{1-\Omega(\tau',T)}{1-\Omega(\tau,T)},\quad
C_\tau = \frac{\mathcal S_P(\kappa,\tau',T)}{\mathcal S_P(\kappa,\tau,T)},
\end{equation}
with $W_\tau$ and $C_\tau$ the weight-effect and composition-effect ratios of
Definition~\ref{def:premium} at this margin. The identity is exact --- it is
\eqref{eq:T1-factor} twice and one division --- and it consumes~(T-1) to guarantee that the two
sides are one model with one primitive changed rather than two different equilibria.

\emph{Step 9 (the weight ratio lies in $[0,1]$).} By leg~1 of Lemma~\ref{lem:L4}, imported by
hypothesis~(T-12) and proved outright there from the clock equivalence of Lemma~\ref{lem:D1}\,(b)
and $b_0 < \tau' < \tau$, we have $\Omega(\tau',T) \ge \Omega(\tau,T)$. Subtracting from one
reverses the inequality, and both sides are strictly positive by~(T-2), so dividing preserves the
direction: $0 \le W_\tau \le 1$, the lower bound because $1-\Omega > 0$ at both policies.

\emph{Step 10 (the composition ratio lies in $[0,1]$).} By leg~3 of Lemma~\ref{lem:L4},
$\mathcal S_P(\kappa,\tau',T) \le \mathcal S_P(\kappa,\tau,T)$, and that leg holds only under
Assumption~\ref{asm:Abr}. With~(T-3) supplying a strictly positive denominator, and with
$\mathcal S_P$ an absolute value by Definition~\ref{def:premium}, $0 \le C_\tau \le 1$. The chain
leg~3 executes is cited rather than re-derived, but it is worth setting out so that each clause can
be seen carrying its own inch of the argument.
\begin{enumerate}[label=(\roman*),leftmargin=2.4em,itemsep=3pt]
\item A tighter threshold weakly lowers the pooled engagement \emph{share}
  $\bar\pi_{\mathrm{pr}} = \Prb(a=1\mid D=0)$. This is leg~2 of Lemma~\ref{lem:L4}, proved outright
  by conditional-probability arithmetic on the nested flagged sets of leg~1.
\item A weakly lower share gives a weakly lower chord endpoint $\bar\pi$. This is clause~(br-iv) of
  Assumption~\ref{asm:Abr}, and it is the clause that keeps the two objects apart: leg~2 moves a
  share, the chord moves an upper support point, and only~(br-iv) links them. Without it leg~2 says
  nothing about $\bar\pi$.
\item A weakly lower $\bar\pi$ gives a weakly smaller $\lvert C_h(\bar\pi)\rvert$. Two clauses do
  different jobs here. The maintained magnitude monotonicity of $\lvert C_h\rvert$ in $\bar\pi$,
  carried by Assumption~\ref{asm:Atau}, orders two values of \emph{one} functional; and
  clause~(br-v) of Assumption~\ref{asm:Abr} is what makes the values at $\tau$ and at $\tau'$ two
  values of one functional rather than of two, since at fixed policies a change in $\tau$ moves
  which histories are pooled, which moves the pooled price, which enters the kernel
  $h = \pi\,p$ through the entry probability \eqref{eq:entry}. The sign half $C_h \le 0$ of
  Assumption~\ref{asm:Atau} is not used at this leg.
\item $\lvert\Akap(\tau')\rvert \le \lvert\Akap(\tau)\rvert$, which is clause~(br-iii) of
  Assumption~\ref{asm:Abr}.
\item Multiplying the two nonnegative factors of the product form \eqref{eq:T1-chord}, using
  items~(iii) and~(iv), gives leg~3.
\end{enumerate}
Items~(iii)--(v) are unavailable unless the representation \eqref{eq:atau} holds at \emph{both}
compared policies, which is clause~(br-i); and item~(v) is unavailable without clause~(br-ii),
because a composition-through-$\kappa$ remainder in $\partial_\kappa M_P$ would destroy the product
form. The threshold leg therefore rests on Assumption~\ref{asm:Atau} at both policies and on all
five clauses~(br-i)--(br-v) of Assumption~\ref{asm:Abr}.

\emph{Step 11 (Part~(B)'s conclusion).} By Steps~9 and~10 both ratios lie in $[0,1]$, so their
product does: $W_\tau C_\tau \le 1\cdot 1 = 1$. Substituting into \eqref{eq:T1-ratio-tau} and
multiplying through by the strictly positive number $\mathcal S(\kappa,\tau,T)$ gives
\eqref{eq:T1B} and, with it,
$\mathcal S(\kappa,\tau',T) \le \mathcal S(\kappa,\tau,T)$. Threshold tightening attenuates at
fixed policies. The structural point is that no dominance condition appears anywhere in Steps~9--11:
at this margin the weight effect and the composition effect push the same way, so the product is
bounded by one without the two being compared in size. The inequality is weak, and it is an
equality when the move from $\tau$ to $\tau'$ reclassifies no history --- when no pair $(j,s)$ has
$\tau' \le B_j(s,H-T) < \tau$ --- since then $\Omega$, $\bar\pi_{\mathrm{pr}}$, $\bar\pi$ and
$\mathcal S_P$ are all unchanged and both ratios equal one. No strict attenuation is claimed.

\medskip
\noindent\emph{Part~(C): the window margin.}

Throughout Part~(C) the threshold $\tau$ is common to the two rules and the comparison is
$T' < T$, so $T'$ is the tighter window.

\emph{Step 12 (the weight ratio is at most one, and this is proved).} Let $j$ be a Voice plan and
$s$ a signal with $D_j(s;\tau,T) = 1$. By the clock equivalence of Lemma~\ref{lem:D1}\,(b),
$B_j(s,H-T) \ge \tau$. Since $T' < T$ gives $H-T' > H-T$, and since $\partial_d B_j \ge 0$ for Voice
plans by clause~(TR-ii) of Assumption~\ref{asm:TR} as it travels with~(T-5), monotonicity of the
stake path in the date gives $B_j(s,H-T') \ge B_j(s,H-T) \ge \tau$; and the engagement flag $a_j=1$
is untouched by the window. Applying the clock equivalence in the other direction at $T'$ yields
$D_j(s;\tau,T') = 1$. By Assumption~\ref{asm:A4} only Voice plans cross the threshold in the core
model, so no other history has to be checked. Hence $\mathcal C_F(\tau,T) \subseteq
\mathcal C_F(\tau,T')$, so $\Omega(\tau,T') \ge \Omega(\tau,T)$, and, exactly as in Step~9 with
\mbox{(T-2)} supplying a strictly positive denominator and~(T-1) holding the plans fixed across the
two windows,
\[
0 \;\le\; W_T \;=\; \frac{1-\Omega(\tau,T')}{1-\Omega(\tau,T)} \;\le\; 1 .
\]
This inequality is a consequence of Lemma~\ref{lem:D1} and the monotone Voice stake path rather
than a hypothesis of the theorem, and it is the bridge from ``shorter window'' to ``higher flagged
weight'' that the legal clock supplies.

\emph{Step 13 (the composition ratio carries no sign).} No hypothesis of this theorem signs $C_T$.
Three reasons block a window analogue of Lemma~\ref{lem:L4}'s third leg, and any one of them
suffices.
\begin{enumerate}[label=(\roman*),leftmargin=2.4em,itemsep=3pt]
\item Assumption~\ref{asm:Abr} is quantified over the threshold pair. Clause~(br-i) asserts the
  representation under $\tau$ and under $\tau'$; clause~(br-iii) compares $\lvert\Akap(\tau')\rvert$
  with $\lvert\Akap(\tau)\rvert$; clause~(br-iv) asserts one endpoint function at $\tau$ and at
  $\tau'$; clause~(br-v) asserts one chord functional at $\tau$ and at $\tau'$. None of them says
  anything about two window environments, so Step~10's chain has no window instance to run on, and
  Lemma~\ref{lem:L4} has no window leg to cite.
\item The window has a channel the threshold does not. At fixed policies, changing $T$ changes the
  filing date $f_j = c_j + T$ and hence the objects clause~(TR-iii) of Assumption~\ref{asm:TR}
  already signs: $T' < T$ implies $B^F(T') \le B^F(T)$ and $Q^F(T') \ge Q^F(T)$. Trade moves from
  the pooled round into the flagged round on histories flagged under \emph{both} windows. A
  threshold change at a fixed window moves no unit of trade across the filing date in this way; it
  relabels which histories file at all.
\item The pooled cell's own information changes. The pooled public history of
  Definition~\ref{def:prices} is $\mathcal H_d^P = (X_0,\dots,X_d;\ \text{flag landed by }d)$, so a
  pooled history carries the event ``no flag has landed by $d$''. Under a tighter window that event
  excludes more crossing dates --- it excludes $c \le d - T'$ rather than only $c \le d - T$ --- so
  the pooled cell's updating changes even holding the set of pooled histories fixed. Clause~(br-ii)
  requires the support points and the kernel not to move with $\kappa$; here the conditioning event
  itself moves with the policy, which is a composition change of the kind~(br-ii) excludes at the
  $\kappa$ margin and says nothing about at the $T$ margin.
\end{enumerate}

\emph{Step 14 (the exact finite equivalence).} Apply \eqref{eq:T1-factor} at $(\tau,T')$ and at
$(\tau,T)$, which~(T-2) permits at both. By~(T-3) and \eqref{eq:T1-factor},
$\mathcal S(\kappa,\tau,T) > 0$, so
\begin{equation}
\label{eq:T1-ratio-T}
\frac{\mathcal S(\kappa,\tau,T')}{\mathcal S(\kappa,\tau,T)}
\;=\; \frac{\bigl(1-\Omega(\tau,T')\bigr)\mathcal S_P(\kappa,\tau,T')}
            {\bigl(1-\Omega(\tau,T)\bigr)\mathcal S_P(\kappa,\tau,T)}
\;=\; W_T\,C_T,
\qquad
C_T = \frac{\mathcal S_P(\kappa,\tau,T')}{\mathcal S_P(\kappa,\tau,T)} .
\end{equation}
Multiplying an inequality by the strictly positive number $\mathcal S(\kappa,\tau,T)$ preserves it
in both directions, so
\begin{equation}
\label{eq:T1-iff}
\mathcal S(\kappa,\tau,T') \;\le\; \mathcal S(\kappa,\tau,T)
\qquad\Longleftrightarrow\qquad
W_T\,C_T \;\le\; 1 ,
\end{equation}
which is the equivalence of Part~(C). Both directions hold, and neither side is a sign claim: the
identity \eqref{eq:T1-ratio-T} is the falsifiable content at this margin, and the equivalence is
empty of economics until $C_T$ is measured. Since $W_T > 0$, \eqref{eq:T1-iff} rearranges to
$C_T \le 1/W_T = \bigl(1-\Omega(\tau,T)\bigr)/\bigl(1-\Omega(\tau,T')\bigr)$: attenuation holds
exactly when the composition effect does not exceed the reciprocal of the weight effect. That is
what ``the weight effect dominates the composition effect'' means here, and it means nothing more.
In particular it does not mean $\lvert 1-W_T\rvert \ge \lvert 1-C_T\rvert$, and it does not mean
$C_T \le 1$.

\emph{Step 15 (the local form).} Adopt the smooth window interpolation of hypothesis~(T-14), which
is consumed here and nowhere else: there is an open interval of window values containing $[T',T]$
and continuously differentiable extensions $r \mapsto \Omega(r)$ and
$r \mapsto \partial_\kappa M_P(r)$ of the model's integer-valued objects, agreeing with them at
$r = -T$ and $r = -T'$, with $\Omega(r) \in (0,1)$ and $\mathcal S_P(r) > 0$ throughout and with
$r \mapsto \Omega(r)$ weakly increasing. Write $r = r_T = -T$ as in Definition~\ref{def:theta}, so
higher $r$ is tighter, and set $r_0 = -T < r_1 = -T'$. On the interpolating interval
$\mathcal S_P(r) > 0$, so $x \mapsto |x|$ is differentiable at $\partial_\kappa M_P(r) \neq 0$ and
$\mathcal S_P$ is continuously differentiable in $r$ by Step~5's second consequence. Differentiating
\eqref{eq:T1-factor} in $r$,
\[
\partial_{r_T}\mathcal S \;=\; -\,\Omega_{r_T}\,\mathcal S_P \;+\; (1-\Omega)\,\partial_{r_T}\mathcal S_P .
\]
The monotonicity clause of~(T-14) gives $\Omega_{r_T} \ge 0$, so the first term is the attenuating
weight effect and the second is the unsigned composition effect. That sign comes from~(T-14) and not
from Step~12: Step~12 compares two \emph{integer} windows and delivers the inequality at the
endpoints only, and nothing outside~(T-14) forbids an interpolant that dips in between. Dividing by
the strictly positive $\mathcal S = (1-\Omega)\mathcal S_P$ and writing, for this proof only,
\begin{equation}
\label{eq:T1-Lambda}
\Lambda(r) \;:=\; \frac{\partial_{r}\mathcal S_P(r)}{\mathcal S_P(r)} \;-\; \frac{\Omega_{r}(r)}{1-\Omega(r)},
\end{equation}
we get $\partial_{r_T}\mathcal S/\mathcal S = \Lambda(r)$, so
$\operatorname{sgn}(\partial_{r_T}\mathcal S) = \operatorname{sgn}(\Lambda)$ and
\begin{equation}
\label{eq:T1-local}
\partial_{r_T}\mathcal S \;\le\; 0
\qquad\Longleftrightarrow\qquad
\frac{\partial_{r_T}\mathcal S_P}{\mathcal S_P} \;\le\; \frac{\Omega_{r_T}}{1-\Omega} ,
\end{equation}
an equivalence at each $r$, with no sign supplied for either side. Display \eqref{eq:T1-local} is
pure algebra --- a division by a strictly positive number --- and holds whether or not the
interpolant is monotone; only the \emph{reading} of the first term as attenuating uses
$\Omega_{r_T} \ge 0$.

\emph{Step 16 (in what sense \eqref{eq:T1C} is the same criterion as \eqref{eq:T1-iff}).} The
statement of Part~(C) says that \eqref{eq:T1C} holds on average along the tightening path,
integrated over $[-T,-T']$, and exactly in the infinitesimal limit, and that it is false read
pointwise. Each of those three clauses is a separate assertion, and each is proved here. Let
$\Lambda$ be as in \eqref{eq:T1-Lambda}, continuous on $[r_0,r_1]$ by~(T-14).

\emph{(i) The finite product criterion is the sign of the integrated local difference.} By~(T-14),~(T-3)
and \eqref{eq:T1-factor}, $\mathcal S(r) > 0$ and is continuously differentiable on $[r_0,r_1]$, so
$\log\mathcal S$ is continuously differentiable there and the fundamental theorem of calculus gives
$\log\bigl(\mathcal S(r_1)/\mathcal S(r_0)\bigr) = \int_{r_0}^{r_1}\partial_r\log\mathcal S(r)\,dr
= \int_{r_0}^{r_1}\Lambda(r)\,dr$, the second equality by Step~15. By \eqref{eq:T1-ratio-T},
$\mathcal S(r_1)/\mathcal S(r_0) = W_T C_T$, so
\begin{equation}
\label{eq:T1-integral}
W_T\,C_T \;=\; \exp\!\left(\int_{r_0}^{r_1}\Lambda(r)\,dr\right),
\qquad\text{hence}\qquad
W_T\,C_T \le 1 \iff \int_{r_0}^{r_1}\Lambda(r)\,dr \le 0 .
\end{equation}
The exponential is strictly increasing, so the equivalence holds in both directions. Written out,
the integrand is the left side of \eqref{eq:T1C} minus its right side: the product criterion is
the local criterion integrated along the tightening path. This is the ``on average'' reading, and
it is exact at finite scale.

\emph{(ii) Pointwise local implies the finite product.} If $\Lambda(r) \le 0$ at every
$r \in [r_0,r_1]$, the integral of a nonpositive continuous function over that interval is
nonpositive, so $W_T C_T \le 1$ by \eqref{eq:T1-integral}. If in addition $\Lambda < 0$ on a
subinterval of positive length, the integral is strictly negative and $W_T C_T < 1$.

\emph{(iii) The finite product implies the local criterion at one point, and no more --- so read
pointwise, \eqref{eq:T1C} is false.} If $W_T C_T \le 1$ then $\int_{r_0}^{r_1}\Lambda \le 0$ by
\eqref{eq:T1-integral}, and since $\Lambda$ is continuous the mean value theorem for integrals
supplies some $r^\star \in (r_0,r_1)$ with
$\Lambda(r^\star) = (r_1-r_0)^{-1}\int_{r_0}^{r_1}\Lambda \le 0$. The local criterion therefore
holds \emph{somewhere} in the interval. It does not hold everywhere in general, and the
counterexample shape is a configuration \eqref{eq:T1C} is not entitled to exclude: a $\Lambda$ that
is positive on $[r_0,\tfrac12(r_0+r_1)]$ and sufficiently negative on $[\tfrac12(r_0+r_1),r_1]$
integrates to a nonpositive number while violating the local criterion on the first half. The
finite comparison from $T$ to $T'$ then reports attenuation even though an intermediate window
tightening amplifies. The implication in this direction is ``at some point'', never ``at every
point''.

\emph{(iv) The two forms coincide exactly in the infinitesimal limit.} Fix $r_0$ and let
$r_1 \downarrow r_0$. By \eqref{eq:T1-ratio-T}, $W_T C_T = \mathcal S(r_1)/\mathcal S(r_0)$ as a
function of $r_1$, and it equals $1$ at $r_1 = r_0$, so
\[
\lim_{r_1\downarrow r_0}\frac{W_TC_T-1}{r_1-r_0}
\;=\; \frac{d}{dr_1}\left.\frac{\mathcal S(r_1)}{\mathcal S(r_0)}\right|_{r_1=r_0}
\;=\; \frac{\partial_r\mathcal S(r_0)}{\mathcal S(r_0)} \;=\; \Lambda(r_0).
\]
Hence if $\Lambda(r_0) < 0$ then $W_TC_T < 1$ for every $r_1 > r_0$ close enough to $r_0$; and if
$W_TC_T \le 1$ for every such $r_1$, the limit is at most zero, so $\Lambda(r_0) \le 0$. The
boundary case $\Lambda(r_0) = 0$ is undetermined at first order, the finite sign there being decided
by higher-order terms, and it is named rather than resolved.

\emph{Step 17 (no unconditional window sign).} Nothing in hypotheses~(T-1)--(T-15) signs $C_T$, by
Step~13, and nothing in them signs $\Lambda$, by Step~15. So nothing in them signs
$\partial_{r_T}\mathcal S$ or $W_TC_T - 1$. Both branches are consistent with every hypothesis
maintained here: $W_TC_T \le 1$ and $W_TC_T > 1$ each require only a value of $C_T$ that the
hypotheses leave free. Part~(C) is an equivalence and an identity, and the sign is an empirical
question about the composition ratio. This is what the statement's closing clause records, and it is
why the model does not deliver window-margin attenuation as a theorem.

Parts~(A), (B) and~(C) are Steps~5--6, Step~11, and Steps~12--17 respectively.
\end{proof}

\medskip
\noindent\emph{Scope of the argument above.}
Hypothesis~(T-15), threshold-side smoothness, is consumed by no step of the proof above. That is the
content of the statement's ``confirmed non-load-bearing'': Part~(B)'s conclusion \eqref{eq:T1B} is a
finite-difference comparison throughout, and Parts~(A) and~(C) never differentiate in $r_\tau$. A
smooth local reading of Part~(B) is available under~(T-15) and would say that
$\partial_{r_\tau}\mathcal S \le 0$ with slack on both sides of zero, which is the same statement as
``no dominance condition needed''; it is not claimed here, and if~(T-15) fails nothing above moves.

Two limits on the proof are worth restating where the argument is, rather than only beside the
statement. First, Part~(B)'s Steps~7 and~10 consume Assumption~\ref{asm:Atau} at both compared
policies and clauses~(br-i)--(br-v) of Assumption~\ref{asm:Abr}; Remark~\ref{rem:AtauRecord} records
that the support condition of $\Atau$ fails at every non-degenerate node of the implemented
calibration, and clause~(br-i) carries that representation at both thresholds, so at that
calibration Part~(B)'s antecedent is not satisfied and Part~(B) says nothing about the implemented
pooled cell. The proof is untouched by that fact and so is the label. Second, the window-margin
evidence recorded at Remark~\ref{rem:T1record} is node evidence at one calibration; it is not a
sign, and Part~(C) claims none.

\begin{proof}[Proof of Proposition~\ref{prop:C1} (C1)]
Four conventions are fixed first, because the statement's hypotheses are stated against them.

\emph{The margin.} By hypothesis~(C-7) the strictness coordinate is the threshold one,
$r = r_\tau = -\tau$ of Definition~\ref{def:theta}, and $\vartheta = (\kappa,r)$.
Definition~\ref{def:primitives} restricts only the window to integers and places no discreteness on
$\tau$, so $r$ is a coordinate on an interval and the derivatives below have a domain to live on.
The window coordinate is excluded: $T$ takes values in $\{1,\dots,H\}$ there, so $\partial_{r_T}$,
$\partial_{\kappa r_T}\Dact$ and $g_{r_T}^{PE}$ have no meaning without an interpolation of the
model in $T$, and nothing local is claimed at that margin.

\emph{The norm.} By hypothesis~(C-1) one norm $\lVert\cdot\rVert$ is fixed on the cutoff space
$\mathbb R^{J-1}$, and every magnitude bar in the $\bar k_x$ and $\bar k_{\kappa r}$ rows of
Definition~\ref{def:theta}, in $L_{\mathcal R}$, and in \eqref{eq:BGE} is read as the member of the
family that norm induces: $\lVert\cdot\rVert$ itself for the vector-valued objects
$\partial_\kappa\Tmap$, $\partial_r\Tmap$ and $\Tmap_{\kappa r}$; the induced operator norm
$\sup_{\lVert w\rVert\le1}\lVert Aw\rVert$ for the matrices $D_k\Tmap$, $\Tmap_{\kappa k}$ and
$\Tmap_{rk}$; $\sup_{\lVert w_1\rVert,\lVert w_2\rVert\le1}\lVert\Tmap_{kk}[w_1,w_2]\rVert$ for the
vector-valued bilinear form $\Tmap_{kk}$; the \emph{dual} norm
$\lVert\phi\rVert_* = \sup_{\lVert w\rVert\le1}\lvert\phi(w)\rvert$ for the covectors $\Delta_k$,
$\Delta_{\kappa k}$ and $\Delta_{kr}$; and
$\sup_{\lVert w_1\rVert,\lVert w_2\rVert\le1}\lvert\Delta_{kk}[w_1,w_2]\rvert$ for the scalar
bilinear form $\Delta_{kk}$. Two consequences are used. The operator norm must be induced, hence
submultiplicative, because Step~1 needs $\lVert A^n\rVert \le \lVert A\rVert^n$; and the pairing
must be the matching one, because pairing a covector with a vector through a magnitude that is not
the dual norm can make \eqref{eq:BGE} \emph{understate} the remainder it is meant to bound, which
voids the implication silently. At $J-1 = 1$ the whole family collapses to the absolute value and
the convention is vacuous.

\emph{The equilibrium objects.} By hypothesis~(C-3) there is, at every $\vartheta$ in the region
$\mathcal R_r$, a cutoff vector $k(\vartheta)$ in the interior of $\Theta$ with
$k(\vartheta) = \Tmap(k(\vartheta);\vartheta)$, and $\vartheta\mapsto k(\vartheta)$ is one branch,
with no switch of selected equilibrium inside $\mathcal R_r$. Interiority is a genuine restriction:
Definition~\ref{def:cpbe}\,(i) has weak inequalities and so permits collapsed action regions, and a
collapsed region puts $k$ on the boundary of $\Theta$, where the fixed-point equation may hold only
as a one-sided condition. Write, for this proof,
\[
\mathcal D(\kappa,r) \;:=\; \Dact\bigl(k(\kappa,r),\kappa,r\bigr),
\qquad
\mathcal S^{GE} \;:=\; \Bigl\lvert \frac{d\mathcal D}{d\kappa} \Bigr\rvert ,
\]
so that $\mathcal D$ is the premium \eqref{eq:Dact} evaluated along the equilibrium branch --- not
the disclosure indicator $D$ --- and $\mathcal S^{GE}$ is the equilibrium liquidity sensitivity of
hypothesis~(C-5). It is a different object from the fixed-policy $\mathcal S = |\partial_\kappa\Dact|$
of Definition~\ref{def:premium}, which holds $k$ frozen; the two agree only where the cutoff
response contributes nothing, and the entire content of this proposition sits in the difference.
Every partial derivative of $\Dact$ and of $\Tmap$ written below is a partial of the
\emph{fixed-policy} map $(k,\kappa,r)\mapsto\Dact(k,\kappa,r)$, respectively of
$(k,\kappa,r)\mapsto\Tmap(k;\kappa,r)$, evaluated at the equilibrium point
$(k(\vartheta),\vartheta)$.

\emph{The region and its regularity.} Hypothesis~(C-1) puts Assumption~\ref{asm:AGE} in force on
$\mathcal R_r$, and names the clause that carries the weight: the contraction bound
$L_{\mathcal R} = \sup_{\mathcal R}\lVert D_k\Tmap\rVert < 1$ of Definition~\ref{def:theta}, read
along the equilibrium path, that is, as the supremum over $\vartheta \in \mathcal R_r$ of
$\lVert D_k\Tmap(k(\vartheta);\vartheta)\rVert$. That reading is the one used at Steps~1--4 and it
is the weaker of the two available; the stronger reading, a supremum over all of $\Theta$ as well,
is not consumed anywhere below. The differentiability clause of Assumption~\ref{asm:AGE} supplies
twice continuous differentiability of $\Tmap$ on a neighbourhood of the equilibrium graph
$\{(k(\vartheta),\vartheta) : \vartheta\in\mathcal R_r\}$, which is the weakest domain Steps~1--4
use; smoothness away from that graph is never called on. Hypothesis~(C-4) supplies twice continuous
differentiability of $\Dact$ in $(k,\kappa,r)$ on the same neighbourhood, and with it finiteness of
the derivatives named in \eqref{eq:BGE} and in $\bar k_{\kappa r}$. The integrability clause of
Assumption~\ref{asm:A2p} is not cited for that and cannot be: boundedness is not differentiability.
Finally, hypothesis~(C-2) makes $\mathcal R_r$ relatively open in both coordinates with
$\kappa\notin\{0,1\}$, so that every $\vartheta\in\mathcal R_r$ has a full two-dimensional
neighbourhood inside $\mathcal R_r$ --- which the implicit function theorem of Step~2, the mixed
partial of Step~5 and the sign argument of Step~7 all need, and which openness in $r$ alone would
not give, $\kappa$ ranging over $[0,1]$ in Definition~\ref{def:primitives}.

\emph{Step 1 (invertibility, without forming an inverse).} Fix $\vartheta\in\mathcal R_r$ and write
$A = D_k\Tmap(k(\vartheta);\vartheta)$, so $\lVert A\rVert \le L_{\mathcal R} < 1$ by~(C-1). By
submultiplicativity, $\lVert A^n\rVert \le L_{\mathcal R}^{\,n}$, so
$\sum_{n\ge0}\lVert A^n\rVert \le (1-L_{\mathcal R})^{-1} < \infty$. The space of
$(J-1)\times(J-1)$ matrices being finite dimensional and complete, the partial sums
$S_N = \sum_{n=0}^{N}A^n$ converge to a limit $S$ with
$\lVert S\rVert \le (1-L_{\mathcal R})^{-1}$; and from $(I-A)S_N = S_N(I-A) = I - A^{N+1}$ with
$\lVert A^{N+1}\rVert \le L_{\mathcal R}^{\,N+1} \to 0$, passing to the limit gives
$(I-A)S = S(I-A) = I$. Hence $I - D_k\Tmap$ is invertible at every point of the equilibrium graph,
with
\begin{equation}
\label{eq:C1-neumann}
\bigl\lVert (I-D_k\Tmap)^{-1} \bigr\rVert \;\le\; \frac{1}{1-L_{\mathcal R}} .
\end{equation}
The bound is a geometric sum of norms, so every derivative bound below is available from
$\lVert D_k\Tmap\rVert$ and one directional derivative of $\Tmap$, with no linear system solved.

\emph{Step 2 (the equilibrium branch is twice continuously differentiable).} Fix
$\vartheta\in\mathcal R_r$ and write $k^0 = k(\vartheta)$. Pick $L'$ with
$L_{\mathcal R} < L' < 1$. By the differentiability clause of Assumption~\ref{asm:AGE} the map
$(k,\vartheta')\mapsto D_k\Tmap(k;\vartheta')$ is continuous, and by~(C-3) and~(C-2) there are
$\delta,\varrho > 0$ with the closed ball $\bar B(k^0,\delta)$ inside the interior of $\Theta$, the
closed ball $\bar B(\vartheta,\varrho)$ inside $\mathcal R_r$ and inside the domain on which
Assumption~\ref{asm:AGE} and~(C-4) hold, and $\lVert D_k\Tmap(k';\vartheta')\rVert \le L'$
throughout $\bar B(k^0,\delta)\times\bar B(\vartheta,\varrho)$. The ball being convex, the
mean-value inequality along the segment gives
$\lVert \Tmap(k_1;\vartheta') - \Tmap(k_2;\vartheta')\rVert \le L'\lVert k_1-k_2\rVert$ there.
Shrinking $\varrho$ so that
$\lVert \Tmap(k^0;\vartheta') - \Tmap(k^0;\vartheta)\rVert \le (1-L')\delta$, and using
$\Tmap(k^0;\vartheta) = k^0$ from~(C-3), gives
$\lVert \Tmap(k_1;\vartheta') - k^0\rVert \le L'\delta + (1-L')\delta = \delta$ for every
$k_1\in\bar B(k^0,\delta)$ and $\vartheta'\in\bar B(\vartheta,\varrho)$. So
$\Tmap(\cdot;\vartheta')$ maps the closed ball into itself and contracts it. That ball is a complete
metric space, so by the contraction mapping theorem each such $\vartheta'$ admits exactly one fixed
point $\hat k(\vartheta')$ of $\Tmap(\cdot;\vartheta')$ inside the ball, with
$\hat k(\vartheta) = k^0$; and the standard estimate
$\lVert\hat k(\vartheta')-\hat k(\vartheta'')\rVert \le (1-L')^{-1}\lVert
\Tmap(\hat k(\vartheta'');\vartheta') - \Tmap(\hat k(\vartheta'');\vartheta'')\rVert$ makes
$\hat k$ continuous. Local uniqueness and local continuity are consequences here, not hypotheses.

Now put $\Phi(k',\vartheta') := k' - \Tmap(k';\vartheta')$ on
$\bar B(k^0,\delta)\times\bar B(\vartheta,\varrho)$. It is twice continuously differentiable by
Assumption~\ref{asm:AGE}, it vanishes at $(k^0,\vartheta)$ by~(C-3), and
$D_k\Phi = I - D_k\Tmap$ is invertible there by Step~1. The implicit function theorem in its twice
continuously differentiable form supplies neighbourhoods and a twice continuously differentiable map
$\vartheta'\mapsto\tilde k(\vartheta')$ solving $\Phi(\tilde k(\vartheta'),\vartheta') = 0$ with
$\tilde k(\vartheta) = k^0$, unique among solutions there; and since every such $\tilde k$ is a
fixed point of $\Tmap(\cdot;\vartheta')$ inside the ball, $\tilde k = \hat k$ where both are
defined. Hypothesis~(C-3)'s single-branch clause is consumed exactly here: the selected
$\vartheta'\mapsto k(\vartheta')$ does not jump to a different fixed point at nearby parameters, so
it stays in $\bar B(k^0,\delta)$ near $\vartheta$ and coincides there with $\hat k = \tilde k$. As
$\vartheta$ was an arbitrary point of the relatively open $\mathcal R_r$, the branch $k(\cdot)$ is
twice continuously differentiable on $\mathcal R_r$.

\emph{Step 3 (the first-derivative bounds).} Let $x\in\{\kappa,r\}$. Differentiating the identity
$k(\vartheta) = \Tmap(k(\vartheta);\vartheta)$ in $x$, which Step~2 licenses and which~(C-2) makes
meaningful in both coordinates, and applying the chain rule in the first argument,
$\partial_x k = D_k\Tmap\,\partial_x k + \partial_x\Tmap$, that is
$(I-D_k\Tmap)\,\partial_x k = \partial_x\Tmap$. Taking norms and using
$\lVert\partial_x k\rVert \le L_{\mathcal R}\lVert\partial_x k\rVert + \lVert\partial_x\Tmap\rVert$,
which may be rearranged because $\lVert\partial_x k\rVert$ is finite by Step~2 and
$L_{\mathcal R} < 1$,
\begin{equation}
\label{eq:C1-kbar}
\lVert\partial_\kappa k\rVert \;\le\; \bar k_\kappa = \frac{\lvert\partial_\kappa\Tmap\rvert}{1-L_{\mathcal R}},
\qquad
\lVert\partial_r k\rVert \;\le\; \bar k_r = \frac{\lvert\partial_r\Tmap\rvert}{1-L_{\mathcal R}},
\end{equation}
which are the $\bar k_x$ rows of Definition~\ref{def:theta}. The derivative on the right is the
\emph{partial} derivative of the outer map in the parameter at frozen cutoffs, not a total
derivative.

\emph{Step 4 (the cross-derivative bound).} Differentiate the $x = \kappa$ instance of Step~3 once
more in $r$, with every $\Tmap$ derivative evaluated at $(k(\vartheta),\vartheta)$ and the chain
rule applied through $k(\vartheta)$ in each slot. In components,
\[
\partial^2_{\kappa r}k^i
= \sum_j\Bigl[\Bigl(\sum_l\Tmap^i_{k_jk_l}\,\partial_r k^l + \Tmap^i_{k_jr}\Bigr)\partial_\kappa k^j
+ \Tmap^i_{k_j}\,\partial^2_{\kappa r}k^j\Bigr]
+ \sum_l\Tmap^i_{\kappa k_l}\,\partial_r k^l + \Tmap^i_{\kappa r},
\]
that is,
$(I-D_k\Tmap)\,\partial^2_{\kappa r}k = \Tmap_{\kappa r} + \Tmap_{\kappa k}[\partial_r k]
+ \Tmap_{rk}[\partial_\kappa k] + \Tmap_{kk}[\partial_r k,\partial_\kappa k]$. Taking norms on the
right with the pairings fixed above and the bounds \eqref{eq:C1-kbar}, then rearranging as in
Step~3,
\begin{equation}
\label{eq:C1-kkr}
\lVert\partial^2_{\kappa r}k\rVert
\;\le\; \bar k_{\kappa r}
= \frac{\lvert\Tmap_{\kappa r}\rvert + \lvert\Tmap_{\kappa k}\rvert\,\bar k_r
+ \lvert\Tmap_{rk}\rvert\,\bar k_\kappa + \lvert\Tmap_{kk}\rvert\,\bar k_\kappa\bar k_r}
{1-L_{\mathcal R}},
\end{equation}
which is the $\bar k_{\kappa r}$ row of Definition~\ref{def:theta}.

\emph{Step 5 (the exact cross-derivative decomposition).} By~(C-4) and Step~2 the composition
$\mathcal D$ is twice continuously differentiable on $\mathcal R_r$, and~(C-2) supplies the
two-dimensional domain the mixed partial needs. The chain rule in $\kappa$ gives
$\partial_\kappa\mathcal D = \Delta_k\cdot\partial_\kappa k + \Delta_\kappa$, and differentiating
that in $r$ term by term, with
$\partial_r(\Delta_k) = \Delta_{kk}[\,\cdot\,,\partial_r k] + \Delta_{kr}$ and
$\partial_r(\Delta_\kappa) = \Delta_{\kappa k}[\partial_r k] + \Delta_{\kappa r}$, gives the exact
identity
\begin{equation}
\label{eq:C1-decomp}
\frac{d^2\mathcal D}{d\kappa\,dr}
\;=\; \underbrace{\Delta_{\kappa r}}_{\text{fixed policy}}
\;+\; \underbrace{\Delta_{\kappa k}[\partial_r k]
\;+\; \bigl(\Delta_{kr} + \Delta_{kk}[\partial_r k]\bigr)\cdot\partial_\kappa k
\;+\; \Delta_k\cdot\partial^2_{\kappa r}k}_{\text{general-equilibrium remainder}} .
\end{equation}
Nothing is dropped and no inequality has been used: \eqref{eq:C1-decomp} is one identity in five
terms, grouped in four. Two features of it are where a miscount would hide. No term
$\Delta_{\kappa\kappa}$ or $\Delta_{rr}$ can appear, the derivative being mixed and each coordinate
differentiated once; and the last term is the only place a \emph{second} derivative of the
equilibrium map enters, which is why Step~4 is needed at all. The remainder is the object
Theorem~\ref{thm:T1} sets its equilibrium free of: the term
$(\partial_\kappa\Omega)(M_F - M_P)$ discarded at \eqref{eq:T1-three} under hypothesis~(T-7) is a
piece of $\Delta_k\cdot\partial_\kappa k$, and it reappears here.

\emph{Step 6 (the remainder is bounded by $\mathcal B_r^{GE}$).} Write $\mathcal E$ for the
remainder group of \eqref{eq:C1-decomp}, so that
$\mathcal E = d^2\mathcal D/d\kappa\,dr - \partial_{\kappa r}\Dact$. Apply the triangle inequality
to its four terms, bound each with the pairings fixed above --- covector against vector through the
dual norm, bilinear form against two vectors --- and insert \eqref{eq:C1-kbar} and
\eqref{eq:C1-kkr}, every factor being finite by~(C-4):
\[
\bigl\lvert\Delta_{\kappa k}[\partial_r k]\bigr\rvert \le \lvert\Delta_{\kappa k}\rvert\bar k_r,
\qquad
\bigl\lvert\bigl(\Delta_{kr}+\Delta_{kk}[\partial_r k]\bigr)\cdot\partial_\kappa k\bigr\rvert
\le \bigl(\lvert\Delta_{kr}\rvert + \lvert\Delta_{kk}\rvert\bar k_r\bigr)\bar k_\kappa,
\qquad
\bigl\lvert\Delta_k\cdot\partial^2_{\kappa r}k\bigr\rvert \le \lvert\Delta_k\rvert\bar k_{\kappa r}.
\]
Summing the three,
\begin{equation}
\label{eq:C1-remainder}
\bigl\lvert \mathcal E \bigr\rvert
\;\le\; \lvert\Delta_{\kappa k}\rvert\bar k_r
+ \bigl(\lvert\Delta_{kr}\rvert + \lvert\Delta_{kk}\rvert\bar k_r\bigr)\bar k_\kappa
+ \lvert\Delta_k\rvert\bar k_{\kappa r}
\;=\; \mathcal B_r^{GE},
\end{equation}
which is \eqref{eq:BGE}. The bookkeeping is worth auditing, since this is where the argument turns:
\eqref{eq:BGE} has three groups and $\mathcal E$ has four terms, mapped one to the first group, two
to the second, and one to the third; every cross-term of the mixed second derivative lands in
exactly one group, no group is fed by a term outside $\mathcal E$, and the only term of
\eqref{eq:C1-decomp} outside the groups is the fixed-policy cross-derivative that
\eqref{eq:C1-remainder} subtracts off. The only inequalities used are the triangle inequality and
the norm pairings. One consequence of \eqref{eq:C1-kbar} and \eqref{eq:C1-kkr} is worth naming: each
$\bar k_x$ carries one factor of $(1-L_{\mathcal R})^{-1}$ and $\bar k_{\kappa r}$ carries three, so
$\mathcal B_r^{GE} = O\bigl((1-L_{\mathcal R})^{-3}\bigr)$ as $L_{\mathcal R}\uparrow1$.

\emph{Step 7 (differentiating the equilibrium sensitivity).} Write $f(\vartheta') :=
d\mathcal D/d\kappa$ at $\vartheta'$, which is continuously differentiable on $\mathcal R_r$ by
Step~5. By hypothesis~(C-5), $f$ does not vanish on $\mathcal R_r$. Take a connected relatively open
neighbourhood $\mathcal N\subseteq\mathcal R_r$ of $\vartheta$, which~(C-2) supplies. A continuous
nowhere-zero real function on a connected set takes values of one sign: otherwise the intermediate
value theorem, applied along a path in $\mathcal N$ joining a point where $f>0$ to one where $f<0$,
would produce a zero of $f$ inside $\mathcal N$, contradicting~(C-5). So
$\operatorname{sgn} f$ is constant on $\mathcal N$; write $\mathfrak s\in\{-1,+1\}$ for that
constant. On $\mathcal N$,
$\mathcal S^{GE} = \lvert f\rvert = \mathfrak s\,f$, an identity of functions and not merely of
values at one point, and the right side is continuously differentiable as a constant multiple of a
continuously differentiable function. Hence $\mathcal S^{GE}$ is differentiable at $\vartheta$ with
\begin{equation}
\label{eq:C1-deriv}
\partial_r\mathcal S^{GE} \;=\; \mathfrak s\,\frac{d^2\mathcal D}{d\kappa\,dr} .
\end{equation}
What~(C-5) does here is remove the single point at which $x\mapsto|x|$ is not differentiable;
without it the conclusion is false, because at a parameter where the equilibrium liquidity
derivative vanishes $\mathcal S^{GE}$ has a kink and $\partial_r\mathcal S^{GE}$ does not exist. The
constancy of the sign is \emph{derived} on $\mathcal N$ from~(C-5) and~(C-2) and is not assumed.

\emph{Step 8 (splitting off the direct margin).} Substitute \eqref{eq:C1-decomp} into
\eqref{eq:C1-deriv} and use $\mathfrak s^2 = 1$. The fixed-policy group is read through
\eqref{eq:gPE}, whose sign factor is the sign of the \emph{equilibrium} liquidity derivative
$d\Dact/d\kappa$ evaluated along $k(\kappa,r)$ --- that is, $\mathfrak s$ --- while its
cross-derivative factor $\partial_{\kappa r}\Dact$ is the fixed-policy partial at
$(k(\vartheta),\vartheta)$. Hence $\mathfrak s\,\Delta_{\kappa r} = -g_r^{PE}$, and
\begin{equation}
\label{eq:C1-split}
\partial_r\mathcal S^{GE} \;=\; -\,g_r^{PE} \;+\; \mathfrak s\,\mathcal E ,
\end{equation}
every object in which is finite by~(C-4).

\emph{Step 9 (dominance closes the argument).} Multiplying by a factor of modulus one leaves the
modulus alone, so $\lvert \mathfrak s\,\mathcal E\rvert = \lvert\mathcal E\rvert \le
\mathcal B_r^{GE}$ by \eqref{eq:C1-remainder}; only the one-sided half
$\mathfrak s\,\mathcal E \le \mathcal B_r^{GE}$ is consumed. With \eqref{eq:C1-split},
\begin{equation}
\label{eq:C1-eta}
\partial_r\mathcal S^{GE}
\;\le\; -\,g_r^{PE} + \mathcal B_r^{GE}
\;=\; -\bigl(g_r^{PE} - \mathcal B_r^{GE}\bigr)
\;=\; -\,\eta_r ,
\end{equation}
with $\eta_r$ the slack of Definition~\ref{def:theta}; and $\eta_r > 0$ at every
$\vartheta\in\mathcal R_r$ by the strict dominance of hypothesis~(C-6). Therefore
$\partial_r\mathcal S^{GE} \le -\eta_r < 0$ at every $\vartheta\in\mathcal R_r$, which is
\eqref{eq:C1}. The argument never needs the remainder's sign, only its size, which is the point of
bounding the general-equilibrium channel rather than signing it.
\end{proof}

\medskip
\noindent\emph{Sign coherence, the norm, and what the argument above does not deliver.}
Three things are worth stating beside the argument rather than only beside the statement.

\emph{Sign coherence is not consumed.} Neither reading of the sign clause reaches \eqref{eq:C1}. The
constancy clause of Assumption~\ref{asm:AGE} --- that the sign of the equilibrium liquidity
derivative is constant on the region --- is not used: Step~7 derives constancy on a connected
neighbourhood from hypotheses~(C-5) and~(C-2), and \eqref{eq:C1-deriv} is a pointwise statement, so
the clause would matter only for naming one sign across a disconnected region. Nor is the coherence
of the fixed-policy sign with the equilibrium sign used. That coherence has one job and it is a job
of naming: at frozen $k$ the fixed-policy sensitivity $\mathcal S$ of Definition~\ref{def:premium}
has $r$-derivative $\operatorname{sgn}(\partial_\kappa\Dact)\,\partial_{\kappa r}\Dact$, and only
when that leading sign equals $\mathfrak s$ does the fixed-policy attenuation rate equal $-g_r^{PE}$
and the phrase ``the fixed-policy attenuation sign survives in equilibrium'' describe
\eqref{eq:C1}. Steps~1--9 are untouched either way, because \eqref{eq:C1-split} pulls out
$\mathfrak s$ and never the fixed-policy sign.

\emph{The hypotheses are relative to the norm.} Every quantity in \eqref{eq:C1-neumann} through
\eqref{eq:C1-eta} other than $\partial_{\kappa r}\Dact$ and $\mathfrak s$ depends on the norm fixed
by hypothesis~(C-1): $L_{\mathcal R}$, the bounds $\bar k_x$ and $\bar k_{\kappa r}$,
$\mathcal B_r^{GE}$, and hence $\eta_r$. A matrix with spectral radius below one has induced
operator norm below one in some norm but not necessarily in a given one, so both the hypothesis
$L_{\mathcal R} < 1$ and the slack in the conclusion are properties of the pair (region, norm), and
a rescaling of the cutoff coordinates can move a parameter vector across the boundary in either
direction. This is not a free lunch, since a change of norm also moves the dual norms inside
$\mathcal B_r^{GE}$; it is a reason to read the hypothesis existentially in the norm and to fix one
convention and keep to it, which is what~(C-1) does.

\emph{What is not delivered.} Nonemptiness of $\mathcal R_r$ is neither claimed nor provable here:
the region is a hypothesis, and whether any parameter vector satisfies (C-1)--(C-7) together is a
property of the calibration. Remark~\ref{rem:C1record} reports the dominance-and-contraction nodes
found by the executed committed check, and what those nodes verify is narrow: the two pointwise
inequalities $L_{\mathcal R} < 1$ and $\eta_r > 0$ together with supporting diagnostics, and not the
full antecedent (C-1)--(C-7) and not a named nonempty region. Promoting nodes to a region needs a
modulus of continuity for $\eta_r$ on the hull of adjacent nodes and a genuine supremum for
$L_{\mathcal R}$ over the region, neither of which more grid nodes can supply. Finally, because
$\mathcal B_r^{GE}$ is a triangle-inequality bound rather than a tight one, $\eta_r \le 0$ carries
no information about the sign of $\partial_r\mathcal S^{GE}$: a failure of dominance is compatible
with equilibrium attenuation, with equilibrium amplification, and with equality.
